% ──────────────────────────────────────────────────────────────────────────────
%  Section 3 — Modulus of the phase current as a real quadratic in (P, Q)
% ──────────────────────────────────────────────────────────────────────────────

\section{\texorpdfstring{Modulus of the phase current as a quadratic form in $(P,Q)$}{Modulus of the phase current as a quadratic form in (P,Q)}}
\label{sec:quadratic}

Taking modulus squared of \autoref{eq:Ik-affine} and expanding,
%
\begin{equation}
\label{eq:Ik-modsq-complex}
\bigl|I_{k}(S,\ell)\bigr|^{2}
\;=\;
\bigl|a_{k}(\ell)\bigr|^{2}\,|S|^{2}
\;+\;
2\,\Re\!\bigl[\,a_{k}(\ell)\,\overline{b_{k}(\ell)}\,S^{*}\,\bigr]
\;+\;
\bigl|b_{k}(\ell)\bigr|^{2}.
\end{equation}
%
Splitting $S = P + \jj Q$ and using
$\Re(zS^{*}) = \Re(z)\,P + \Im(z)\,Q$ gives the equivalent real-valued
form
%
\begin{equation}
\label{eq:Ik-modsq-real}
\bigl|I_{k}(P,Q,\ell)\bigr|^{2}
\;=\;
\xi_{k}(\ell)\,(P^{2}+Q^{2})
\;+\;
2\,\Re\!\bigl[\eta_{k}(\ell)\bigr]\,P
\;+\;
2\,\Im\!\bigl[\eta_{k}(\ell)\bigr]\,Q
\;+\;
\zeta_{k}(\ell),
\end{equation}
%
\noindent where the three line-only scalars per phase are
%
\begin{equation}
\label{eq:xi-eta-zeta}
\xi_{k}(\ell) \;=\; \bigl|a_{k}(\ell)\bigr|^{2},
\qquad
\eta_{k}(\ell) \;=\; a_{k}(\ell)\,\overline{b_{k}(\ell)},
\qquad
\zeta_{k}(\ell) \;=\; \bigl|b_{k}(\ell)\bigr|^{2}.
\end{equation}
%
Among these, $\xi_{k}$ and $\zeta_{k}$ are non-negative reals and
$\eta_{k}$ is complex. They are derived once from \autoref{eq:ak-bk}
for a given length.

\autoref{eq:Ik-modsq-real} states that the modulus squared of every
phase current is a quadratic polynomial in the active and reactive
power, with no cross product $PQ$, with the coefficient of
$(P^{2}+Q^{2})$ equal across the three phases up to $\xi_{k}$, and with
the linear part fully captured by the two real coordinates of the
complex constant $\eta_{k}$. The asymmetry of the line lives entirely
in the spread of $\{\xi_{k},\,\eta_{k},\,\zeta_{k}\}$ across
$k\in\{A,B,C\}$.

The phase current modulus itself is the square root:
%
\begin{equation}
\label{eq:Ik-mod}
\bigl|I_{k}(P,Q,\ell)\bigr|
\;=\;
\sqrt{\,
\xi_{k}(\ell)\,(P^{2}+Q^{2})
\;+\;
2\,\Re[\eta_{k}(\ell)]\,P
\;+\;
2\,\Im[\eta_{k}(\ell)]\,Q
\;+\;
\zeta_{k}(\ell)\,}.
\end{equation}
%
For fixed $\ell$ the level sets of \autoref{eq:Ik-mod} in the $(P,Q)$
plane are circles centred at
$\bigl(-\Re[\eta_{k}]/\xi_{k},\,-\Im[\eta_{k}]/\xi_{k}\bigr)$;
the displacement of those three centres is one geometric reading of the
phase-current unbalance.
